\documentclass[11pt]{article}
\usepackage{amsmath,amssymb,amsthm}
\usepackage[margin=1.15in]{geometry}
\usepackage{microtype}

\theoremstyle{plain}
\newtheorem{theorem}{Theorem}
\newtheorem{proposition}[theorem]{Proposition}
\newtheorem{corollary}[theorem]{Corollary}
\theoremstyle{definition}
\newtheorem{definition}[theorem]{Definition}
\newtheorem{example}[theorem]{Example}
\theoremstyle{remark}
\newtheorem{remark}[theorem]{Remark}

\newcommand{\ist}{\mathrm{ist}}
\newcommand{\Cn}{\mathrm{Cn}}
\newcommand{\Aut}{\mathrm{Aut}}
\newcommand{\out}{\mathrm{out}}
\newcommand{\Lp}{L^{+}}

\title{Uniformity and Transcendence\\
\large A trichotomy for lifting schemas in the logic of context}
\author{}
\date{}

\begin{document}
\maketitle

\begin{abstract}
McCarthy's programme for a logic of context asks for \emph{transcendence}: the
ability to step outside any context and thereby reach conclusions unavailable
within it. We show that transcendence is incompatible with uniformity, in a
precise sense. A lifting schema that treats all propositions alike either
yields nothing new at the base level, or fails to determine a unique extension,
or is not in fact a schema but a disguised assertion. In the monotone case the
trichotomy collapses to a dichotomy: monotone uniform schemas are always
conservative. We further show that a schema uniform in the \emph{contexts} as
well has a shift-invariant extension --- every level of the tower says exactly
what every other level says --- which is McCarthy's ``pointless version''
derived rather than stipulated. All three results identify uniformity, not
consistency, as the binding constraint. The harmlessness of the regress is
thereby explained: it is harmless because it is inert, and it is inert because
it is uniform. We intend this as an explanation of why the transcendence
programme stalled rather than a refutation of it, and we draw from it a positive
design principle: content in a context system comes from deliberate asymmetry,
never from the tower. Cashed out, that principle is an audit property --- a
record whose update rule is uniform says exactly what was logged, and every fact
in it is attributable to a revision.
\end{abstract}

\section{Introduction}

The infinite tower of contexts $c_0, c_1, c_2, \dots$, in which each context is
asserted inside the next, is consistent. That much is settled: the tower is
satisfiable, and satisfiable transfinitely, by an evaluator respecting the
connectives. The regress is harmless.

This paper takes up the question that harmlessness leaves open. McCarthy did not
want the tower merely to be safe; he wanted it to be \emph{productive}. Entering
an outer context was supposed to license conclusions the inner one could not
reach. Our results say this cannot happen, and they locate the obstruction
precisely --- not in consistency, and not in the infinitude of the regress, but
in the demand that a lifting rule be a \emph{rule}.

Nothing below contradicts anything McCarthy asserted. The tower is consistent,
the machinery of $\ist$ and lifting is sound, and the transcendence programme was
always presented as preliminary and unfinished rather than as a completed
theory. Our results supply the missing account of \emph{why} it remained
unfinished. The analogy is with a decision problem answered in the negative: the
question was worth asking, the answer is no, and the proof that it is no is
itself informative about the subject.

The informal claim we make rigorous is:

\begin{quote}
You may have new conclusions, or you may have a determinate answer, or you may
have a uniform rule. You may not have all three.
\end{quote}

Section 2 fixes the framework. Section 3 proves the monotone dichotomy
(Theorem \ref{thm:mono}). Section 4 proves the general trichotomy
(Theorem \ref{thm:trichotomy}). Section 5 proves the rigidity theorem
(Theorem \ref{thm:rigidity}), which recovers the copy schema from
context-uniformity alone. Section 6 turns the negative results into a positive
one, about records that must be auditable. Section 7 discusses the hypotheses,
and is where a sceptical reader should begin.

\section{Framework}

\subsection{Languages and frames}

Let $L$ be a propositional language over a countable set $A$ of atoms, closed
under the Boolean connectives. $L$ is the \emph{base language}: the language in
which facts about the world are stated, containing no context machinery. We
write $\Cn$ for classical consequence.

\begin{definition}[Frame]
A \emph{frame} is a pair $F=(C,\out)$ where $C$ is a nonempty set of contexts and
$\out\colon C \to C$ is a total injection. Totality is McCarthy's requirement
that every context lie inside a further context; injectivity says a context has
at most one immediate interior.
\end{definition}

Two frames matter below. The \emph{tower frame} $\mathbb{N}$ has
$C=\{c_n : n \in \mathbb{N}\}$ and $\out(c_n)=c_{n+1}$; here $c_0$ is the unique
non-successor. The \emph{bi-infinite frame} $\mathbb{Z}$ has $C=\{c_n : n \in
\mathbb{Z}\}$ and the same map; here $\out$ is a bijection and no context is
distinguished.

\begin{definition}[Context language]
$\Lp(F)$ is the least set containing $A$, closed under the Boolean connectives,
and such that $\ist_c(\varphi) \in \Lp(F)$ whenever $c \in C$ and
$\varphi \in \Lp(F)$.
\end{definition}

\begin{definition}[Shift]
The \emph{shift} $\sigma_F \colon \Lp(F)\to\Lp(F)$ is defined by
$\sigma_F(p)=p$ for $p \in A$, commutation with the connectives, and
$\sigma_F(\ist_c(\varphi)) = \ist_{\out(c)}(\sigma_F \varphi)$.
\end{definition}

On $\mathbb{N}$ the shift is an injection whose image omits every formula whose
outermost modality is $\ist_{c_0}$. On $\mathbb{Z}$ it is a bijection. This
asymmetry is the whole content of Section 5.

\subsection{Schemas and extensions}

We allow lifting rules to be defeasible, since the interesting candidates for
transcendence are.

\begin{definition}[Schema]
A \emph{schema} $S$ assigns to each frame $F$ a set $S_F$ of default rules
$\dfrac{\alpha : \beta}{\gamma}$ over $\Lp(F)$, with $\alpha$ the prerequisite,
$\beta$ the justification and $\gamma$ the consequent. A rule with
$\beta = \top$ is \emph{monotone}; $S$ is monotone if all its rules are.
\end{definition}

\begin{definition}[Extension]\label{def:extension}
Given a theory $T$ and a schema $S$ on $F$, a deductively closed
$E \subseteq \Lp(F)$ is an \emph{extension} of $(T,S_F)$ if $E = \bigcup_{i}E_i$
where $E_0 = \Cn(T)$ and
\[
E_{i+1} \;=\; \Cn(E_i) \;\cup\;
\Bigl\{\gamma \;:\; \tfrac{\alpha\,:\,\beta}{\gamma}\in S_F,\;
\alpha \in E_i,\; \neg\beta \notin E \Bigr\}.
\]
We count only consistent extensions. If $S$ is monotone this fixed point is
unique; it equals $\Cn(T\cup\Gamma)$, with $\Gamma$ the set of consequents,
when every rule is \emph{prerequisite-free} ($\alpha=\top$), but not in
general --- see Remark~\ref{rem:prereq}.
\end{definition}

\begin{remark}[Prerequisites]\label{rem:prereq}
A monotone rule still carries a prerequisite, and a rule whose prerequisite is
never derived never fires. Over the tautological base the single rule
$\frac{a\,:\,\top}{a}$ has extension $\Cn(\emptyset)$, which does not contain
$a$, whereas $\Cn(\emptyset\cup\{a\})$ does. So the extension of a monotone
schema is in general a proper subset of $\Cn(T\cup\Gamma)$.

What does hold, for \emph{every} schema and with no hypothesis at all, is the
inclusion $E\subseteq\Cn(T\cup\Gamma)$: the base stage is contained in it, and
each successor stage is too, since $\Cn$ is idempotent and every fired
consequent lies in $\Gamma$. That inclusion is all of Section~\ref{sec:records}
needs, which is why nothing there assumes monotonicity.
\end{remark}

\begin{definition}[Conservativity]
$S$ is \emph{conservative} over $T$ if $E \cap L = \Cn(T)\cap L$ for every
extension $E$. \emph{Transcendence} is the failure of conservativity: some
sentence of the base language becomes derivable that was not derivable before.
\end{definition}

We stress that conservativity is measured in $L$, not in $\Lp$. Every nontrivial
schema adds sentences of $\Lp$ --- that is what a schema is for. The question
McCarthy poses is whether climbing the tower tells you anything about the
\emph{world}, and the world is described in $L$.

\subsection{Uniformity}

Two independent ways a purported rule can smuggle in an assertion.

\begin{definition}[Proposition-uniformity]
$S$ is \emph{proposition-uniform} if $s(S_F) \subseteq S_F$ for every
substitution $s \colon A \to L$, extended homomorphically to $\Lp(F)$ (acting
inside the scope of $\ist$).
\end{definition}

A lifting rule should be a rule about lifting, uniform in what is lifted. A
``schema'' that fires only for one designated sentence $\theta$ is an assertion
of $\theta$ wearing a rule's clothing.

\begin{definition}[Context-uniformity]
$S$ is \emph{context-uniform} if $\sigma_F^{-1}(S_F) = S_F$ for every frame $F$,
and $S_F$ is defined for every frame.
\end{definition}

The quantification over \emph{all} frames is essential and is not a technical
convenience; see Remark \ref{rem:definable}.

\begin{definition}
$S$ is \emph{uniform} if it is both proposition- and context-uniform.
\end{definition}

\subsection{The base theory}

Throughout, $T \subseteq L$: the base theory is context-free. This is not
cosmetic. If $T$ may itself contain $\ist$-sentences, it can supply a premise
that a lifting rule then exports, and non-conservativity follows for a reason
having nothing to do with the tower. See Remark \ref{rem:contextfree}.

For Theorem \ref{thm:trichotomy} we need a symmetry condition on $T$.

\begin{definition}[Atom automorphisms]\label{def:aut}
For a permutation $\pi$ of $A$ and an \emph{arbitrary} function
$\varepsilon \colon A \to \{0,1\}$, let $g_{\pi,\varepsilon}$ be the
substitution sending the atom $a$ to $\pi(a)$ if $\varepsilon(a)=0$ and to
$\neg\pi(a)$ if $\varepsilon(a)=1$, extended homomorphically. Each
$g_{\pi,\varepsilon}$ is an automorphism of $L$, and
\[
\Aut(L) \;:=\; \{\, g_{\pi,\varepsilon} \;:\; \pi \in \mathrm{Sym}(A),\;
\varepsilon \in \{0,1\}^{A} \,\}
\;\cong\; \{0,1\}^{A} \rtimes \mathrm{Sym}(A)
\]
is a group under composition. We write $\Aut(L)_T$ for its subgroup preserving
$\Cn(T)$. That $\varepsilon$ ranges over \emph{all} functions, not merely those
of finite support, is essential; see Remark~\ref{rem:finitegen}.
\end{definition}

These are automorphisms of the Lindenbaum algebra of $L$, not of the algebra of
formulas: $g_{\pi,\varepsilon}$ applied twice to the atom $a$ returns
$\neg\neg a$, which is equivalent to $a$ but is a different formula. Nothing in
the proofs below depends on the distinction, since every set they act on is
deductively closed. It matters only if one wants $\Aut(L)$ to act on syntax, for
which it suffices to let atoms carry a sign and read negation as a sign flip;
this is what the machine-checked development does, for the reason given in
Section~\ref{sec:formal}.

Automorphisms act on formulas covariantly and hence on valuations
contravariantly, so we use the induced left action
$g \cdot M := (g^{-1})^{*}M$, where the pullback $h^{*}M$ is the valuation
$a \mapsto M(h(a))$. It is characterised by
\[
g \cdot M \models g\varphi \quad\Longleftrightarrow\quad M \models \varphi ,
\]
so $\Aut(L)_T$ permutes the models of $T$.

\begin{definition}[Symmetric theory]
$T$ is \emph{symmetric} if $\Aut(L)_T$ acts transitively on the models of $T$.
\end{definition}

\begin{proposition}\label{prop:tautsym}
The tautological theory is symmetric.
\end{proposition}

\begin{proof}
Every valuation is a model and $\Aut(L)_{\top}=\Aut(L)$. Given valuations $M,N$,
take $\pi=\mathrm{id}$ and $\varepsilon(a)=1$ exactly when $M(a)\neq N(a)$. Then
$g_{\pi,\varepsilon}$ is an involution and $g_{\pi,\varepsilon}\cdot M = N$.
\end{proof}

This is the case of interest: McCarthy's question is whether the tower is by
itself a source of new facts, so the base should be assumed to contribute
nothing.

\begin{remark}[Why generation is the wrong definition]\label{rem:finitegen}
One is tempted to define $\Aut(L)$ as the group \emph{generated by} permutations
of $A$ together with negations of individual atoms. That group is strictly
smaller. Conjugating a single-atom negation by a permutation yields another
single-atom negation, so the negations generate a normal subgroup consisting
exactly of the $g_{\mathrm{id},\varepsilon}$ with $\varepsilon$ of finite
support, and every element of the generated group is $g_{\pi,\varepsilon}$ with
$\varepsilon$ finitely supported. For infinite $A$ such a group is not
transitive on valuations: the constant-true and constant-false valuations differ
at infinitely many atoms, and no permutation alters that. The tautological
theory would then fail to be symmetric, and Theorem~\ref{thm:trichotomy} would
lose the very instance it is meant to cover. Definition~\ref{def:aut} takes the
full group instead, at no cost to the proof, which uses only transitivity and
preservation of $\Cn(T)$.
\end{remark}

\section{The monotone case}

Here uniqueness is automatic, so the trichotomy has no room to be a trichotomy.
It collapses, and the collapse is sharp.

\begin{theorem}[Monotone dichotomy]\label{thm:mono}
Let $T \subseteq L$ be consistent and let $S$ be monotone and
proposition-uniform. Then either $S_F$ is inconsistent, or $S$ is conservative
over $T$.
\end{theorem}

\begin{proof}
Let $E$ be the unique extension and suppose $\varphi \in E \cap L$ with
$\varphi \notin \Cn(T)$. Fix a model $M \models T$ with $M \not\models \varphi$,
and let $s$ be the substitution sending each atom $p$ to $\top$ if $M \models p$
and to $\bot$ otherwise.

Since $E = \Cn(T \cup \Gamma)$ for $\Gamma$ the set of consequents of $S_F$,
there are $\psi_1,\dots,\psi_m \in T$ and $\gamma_1,\dots,\gamma_k \in \Gamma$
with
\[
\psi_1,\dots,\psi_m,\gamma_1,\dots,\gamma_k \vdash \varphi .
\]
Substitution preserves derivability, so
\[
s\psi_1,\dots,s\psi_m,\,s\gamma_1,\dots,s\gamma_k \vdash s\varphi .
\]
Each $\psi_i \in T \subseteq L$ holds in $M$, so $s\psi_i$ is a closed Boolean
combination of $\top,\bot$ with value $\top$; each is a tautology. Dually
$M \not\models \varphi$ gives $s\varphi \equiv \bot$. By proposition-uniformity
each $s\gamma_j$ is again a consequent of $S_F$. Hence
$\Gamma \vdash \bot$, so $S_F$ is inconsistent.
\end{proof}

\begin{corollary}\label{cor:mono}
A consistent monotone uniform schema cannot deliver transcendence. Since it also
has exactly one extension, uniqueness and novelty are outright incompatible in
the monotone setting.
\end{corollary}

\begin{remark}
Context-uniformity is not used in Theorem \ref{thm:mono}. Only the treatment of
propositions matters. This is the first indication that the level dimension is
not where the obstruction lives.
\end{remark}

\section{The general case}

\begin{theorem}[Trichotomy]\label{thm:trichotomy}
Let $T \subseteq L$ be consistent and symmetric, let $F$ be any frame, and let
$S$ be any schema. Then at least one of the following holds.
\begin{itemize}
\item[\textup{(C)}] $S$ is conservative over $T$.
\item[\textup{(U)}] $(T,S_F)$ has a number of consistent extensions different
from one --- zero, or more than one.
\item[\textup{(P)}] $S$ is not proposition-uniform.
\end{itemize}
\end{theorem}

\begin{proof}
Assume (U) and (P) fail; we derive (C). Let $E$ be the unique consistent
extension.

\emph{Step 1: $E$ is invariant.}
Let $g \in \Aut(L)_T$, extended to $\Lp(F)$ by acting inside the scope of each
$\ist_c$ and fixing every context index. Because $g$ is invertible and $S$ is
proposition-uniform we have $g(S_F) \subseteq S_F$ and
$g^{-1}(S_F)\subseteq S_F$, hence $g(S_F)=S_F$. Also $g(\Cn(T))=\Cn(T)$ by
choice of $g$. Since $g$ is an automorphism of $\Lp(F)$ it commutes with $\Cn$
and preserves consistency, so it carries the fixed-point construction to itself:
$g(E)$ is a consistent extension. By uniqueness $g(E)=E$.

\emph{Step 2: invariance forces conservativity.}
Suppose $\varphi \in E \cap L$ with $\varphi \notin \Cn(T)$, and fix
$M \models T$ with $M \not\models \varphi$. Let $N$ be any model of $T$. By
symmetry there is $g \in \Aut(L)_T$ with $gM = N$, and then
$N \not\models g\varphi$, since $g$ is an isomorphism of the satisfaction
relation. By Step 1, $g\varphi \in E$.

Thus $\{g\varphi : g \in \Aut(L)_T\} \subseteq E \cap L$, and this set is
falsified at every model of $T$. Hence $T \cup \{g\varphi\}_g$ is unsatisfiable,
so $E \supseteq \Cn(T \cup \{g\varphi\}_g)$ is inconsistent, contradicting our
choice of $E$. Therefore $E \cap L = \Cn(T)\cap L$, which is (C).
\end{proof}

\begin{remark}
The three corners are jointly exhaustive and individually occupied. The copy
schema occupies (C). A lifting rule with two conflicting destinations and no
priority between them occupies (U). Example \ref{ex:theta} occupies (P).
\end{remark}

\begin{example}[The disguised assertion]\label{ex:theta}
Fix $\theta \in L$ with $\theta \notin \Cn(T)$ and let
$S_F = \{\ist_c(\theta) : c \in C\}$ together with the copy and an exit rule.
This is monotone, hence has exactly one extension; it is shift-invariant, hence
context-uniform; and it yields $\theta$. It is not proposition-uniform, and this
is the only hypothesis it violates. It shows that (P) cannot be weakened to
context-uniformity alone, and it is the counterexample a referee will construct
first.
\end{example}

\section{Rigidity: recovering the copy}

Theorems \ref{thm:mono} and \ref{thm:trichotomy} say nothing new arrives at the
base. A stronger structural fact holds once contexts are treated uniformly: the
levels become indistinguishable from one another.

\begin{theorem}[Rigidity]\label{thm:rigidity}
Let $T \subseteq L$ be consistent and let $S$ be context-uniform with a unique
consistent extension $E$ on $F$. Then $\sigma_F^{-1}(E) = E$. Consequently, for
all $\varphi \in \Lp(F)$ and all $c$,
\[
\ist_c(\varphi) \in E \quad\Longleftrightarrow\quad
\ist_{\out(c)}(\sigma_F\varphi) \in E ,
\]
and in particular, for $\varphi \in L$, where $\sigma_F$ acts as the identity,
\[
\ist_c(\varphi) \in E \quad\Longleftrightarrow\quad
\ist_{\out(c)}(\varphi) \in E .
\]
Every level of the tower therefore carries the same content read through the
shift, and literally the same content about the base language.
\end{theorem}

\begin{proof}
Write $\sigma=\sigma_F$. We claim $\sigma^{-1}(E)$ is an extension of
$(T,S_F)$.

It contains $T$, since $\sigma$ fixes $L$ pointwise and $T\subseteq L$. It is
deductively closed: $\sigma$ is an injective homomorphism commuting with the
connectives, so if $\sigma^{-1}(E) \vdash \chi$ then $E \vdash \sigma\chi$,
whence $\sigma\chi \in E$ and $\chi \in \sigma^{-1}(E)$. It is closed under
$S_F$: if $\frac{\alpha:\beta}{\gamma} \in S_F$ with
$\alpha \in \sigma^{-1}(E)$ and $\neg\beta \notin \sigma^{-1}(E)$, then by
context-uniformity $\frac{\sigma\alpha:\sigma\beta}{\sigma\gamma} \in S_F$, and
$\sigma\alpha \in E$, $\neg\sigma\beta \notin E$; so $\sigma\gamma \in E$ and
$\gamma \in \sigma^{-1}(E)$. Groundedness is inherited stagewise by the same
argument.

It is consistent, since $\bot \in \sigma^{-1}(E)$ would give
$\bot = \sigma\bot \in E$.

By uniqueness $\sigma^{-1}(E) = E$; that is, $\chi \in E$ iff $\sigma\chi \in E$.
Taking $\chi = \ist_c(\varphi)$ and using
$\sigma(\ist_c(\varphi)) = \ist_{\out(c)}(\sigma\varphi)$ gives the first
displayed equivalence, and $\sigma\!\restriction\! L = \mathrm{id}$ gives the
second.
\end{proof}

\begin{remark}[The shift cannot be dropped]
In the first display the right-hand side must be $\sigma_F\varphi$, not
$\varphi$. The shift relabels every context index occurring inside $\varphi$, so
$\ist_c(\varphi)$ and $\ist_{\out(c)}(\varphi)$ are $\sigma_F$-related only when
$\varphi$ contains no occurrence of $\ist$ --- that is, only on the base
language. For $\varphi = \ist_{c}(p)$ the unshifted equivalence would assert a
coincidence between two unrelated pairs of levels, and nothing in the hypotheses
supplies it.
\end{remark}

\begin{corollary}
Under the hypotheses of Theorems \ref{thm:trichotomy} and \ref{thm:rigidity}
together, the unique extension is the copy: it adds no base-level fact and no
level-to-level variation. This is McCarthy's ``pointless version,'' obtained as
a conclusion rather than adopted as a stipulation.
\end{corollary}

\begin{remark}[The definable-$c_0$ loophole]\label{rem:definable}
On the tower frame $\mathbb{N}$ the context $c_0$ is definable as the unique
non-successor, so a schema may name it without using a constant --- for
instance, ``at the innermost context, the set of truths is complete.'' Such a
schema satisfies $\sigma^{-1}(S)=S$ on $\mathbb{N}$ vacuously at the relevant
place and would evade a frame-local definition of context-uniformity.

Quantifying over all frames closes this. On $\mathbb{Z}$ every context is a
successor, the definition ``innermost'' has empty extension, and the schema goes
silent. A schema whose content depends on there being a bottom is relying on
structure the theory never promised: McCarthy's own axiom is that every context
has an outer one, never that some context has no inner one. The $\mathbb{Z}$
frame is not a trick; it is the theory taken at its word.

Note also that on $\mathbb{Z}$ the shift is a bijection, so
Theorem \ref{thm:rigidity} yields invariance under a genuine automorphism rather
than a mere embedding.
\end{remark}

\section{Records: the principle cashed out}\label{sec:records}

The results so far are limitative. They have a positive corollary, and it is
worth stating as a theorem rather than as advice: if no content-blind rule can
generate content, then every fact a context system delivers is traceable to
something a person put there deliberately. The natural object to say this about
is a record that is revised over time and has to be auditable.

Read the contexts as revisions. At revision $n$ a set $\Lambda_n\subseteq L$ of
\emph{entries} is logged. Entries lie in the base language because they are
claims about the world, not about the record --- the same context-freeness
Theorem~\ref{thm:mono} requires, arrived at for an independent reason. Write
$\Lambda_{\le n}=\bigcup_{i\le n}\Lambda_i$, fix an update schema $S$, and let
$E_n$ be an extension of $(\Lambda_{\le n},S_F)$.

\begin{theorem}[Audit]\label{thm:audit}
Let $S$ be proposition-uniform and let $\Lambda_{\le n}\cup\Gamma$ be
consistent, where $\Gamma$ is the set of consequents of $S_F$. Then for every
$\varphi\in L$,
\[
\varphi\in E_n \quad\Longleftrightarrow\quad \varphi\in\Cn(\Lambda_{\le n}).
\]
\end{theorem}

\begin{proof}
Right to left, $E_n\supseteq\Cn(\Lambda_{\le n})$ holds by construction, since
$E_0=\Cn(\Lambda_{\le n})$ is the base stage. Left to right, every stage lies
inside $\Cn(\Lambda_{\le n}\cup\Gamma)$ by Remark~\ref{rem:prereq}, so
$E_n\subseteq\Cn(\Lambda_{\le n}\cup\Gamma)$; Theorem~\ref{thm:mono}, applied
with base theory $\Lambda_{\le n}$ and consequents $\Gamma$, returns
$\varphi\in\Cn(\Lambda_{\le n})$.
\end{proof}

\begin{corollary}[Attribution]\label{cor:attribution}
If $\varphi\in L$ lies in $E_n$ but not in $E_m$, then
$\varphi\in\Cn(\Lambda_{\le n})$ and $\varphi\notin\Cn(\Lambda_{\le m})$. Some
entry logged after revision $m$ is therefore necessary for $\varphi$.
\end{corollary}

\begin{corollary}[No silent loss]\label{cor:noloss}
For $m\le n$, every base sentence of $E_m$ lies in $E_n$.
\end{corollary}

Together: the record says exactly what was logged, it never says less than it
did, and nothing appears in it that no revision accounts for. That is the
property an auditor is actually asking after, and it is now a consequence of
the monotone dichotomy rather than a design aspiration.

\begin{remark}[Why not the trichotomy]
Theorem~\ref{thm:trichotomy} cannot be used here. It requires the base theory
to be symmetric, and a log is the opposite of symmetric --- it is a specific
claim about the world, which is the whole point of keeping one.
Theorem~\ref{thm:mono} asks only for context-freeness, which a log satisfies by
construction. The negative result with the weaker hypothesis is the one that
turns out to be useful.
\end{remark}

\begin{remark}[The schema may be defeasible]
Theorem~\ref{thm:audit} does not assume $S$ is monotone. Only the inclusion of
Remark~\ref{rem:prereq} is used, and that holds for any schema. So the
practically interesting update rules --- carry the previous value forward unless
something supersedes it --- are covered.

Non-uniqueness is not an obstacle either, merely a different fact. If
$(\Lambda_{\le n},S_F)$ has several consistent extensions, the theorem holds of
each: the record could have come out differently, but whichever way it came out,
it says nothing about the world beyond the log. What multiple extensions cost
is reproducibility, not auditability.
\end{remark}

\begin{remark}[The hypotheses are tight]\label{rem:tight}
Each hypothesis of Theorem~\ref{thm:audit} has a counterexample on the other
side of it, and each corresponds to a recognisable way of building a record
system wrongly.

Drop context-freeness of the entries and the theorem fails by
Remark~\ref{rem:contextfree}: a log that already speaks of revisions can supply
a premise the schema exports.

Drop proposition-uniformity and it fails by Example~\ref{ex:theta}: the schema
is then an assertion in disguise, and it is the assertion, not the record, that
is doing the work.

Add a completeness claim --- an update rule saying that revision $n$ contains
everything true at $n-1$ \emph{and nothing else} --- and it fails for a reason
established in the companion development, where reifying the set of truths is
shown to be conservative but asserting that set complete is shown not to be.
This is the closed-world cache: treating the absence of an entry as the presence
of its negation.

Drop consistency of $\Lambda_{\le n}\cup\Gamma$ and everything follows, which is
not a subtlety but is worth naming: a log that contradicts what the update rule
could ever say is a record that proves anything.
\end{remark}

\section{Discussion}

\subsection{What the programme got right}

It would be a misreading to take the above as a refutation. Three points deserve
to be stated plainly.

First, the impossibility is an impossibility for \emph{uniform} rules, and the
uniformity conditions are ones the programme itself would have insisted on. A
lifting axiom that fires only for a designated sentence is not the kind of thing
McCarthy was looking for. The theorems therefore respect the programme's own
standards; they show that those standards, honestly applied, exclude the
outcome.

Second, the copy schema deserves better than the label ``pointless.'' Judged as
a source of new theorems it is empty, and that judgment is correct. Judged as a
guarantee of stability it is exactly what one wants: a system whose successive
contexts are faithful copies of their predecessors cannot silently revise itself.
Where the requirement is the absence of drift rather than the production of
novelty --- persistent records, audit trails, any setting in which a commitment
made earlier must still mean the same thing later --- inertness is the
deliverable and not the defect. The copy schema is the formal content of that
guarantee, and Theorem \ref{thm:rigidity} shows it is not merely one option among
many but the forced consequence of treating contexts uniformly.

Third, and most usefully, the trichotomy converts a negative result into a
design principle. If uniform rules cannot generate content, then every genuine
inference a context system performs is traceable to an asymmetry someone put
there on purpose: a domain axiom, a priority ordering, a distinguished context, a
base theory that is not symmetric. This is a sharper piece of guidance than the
transcendence hope ever provided, because it tells the designer exactly where to
look when a system draws a conclusion nobody intended. The answer is never the
tower.

\subsection{What is doing the work}

The obstruction is not consistency. The tower is consistent, at every finite and
transfinite stage, and no result above uses the failure of consistency anywhere.
It is not the infinitude of the regress either: Theorem \ref{thm:mono} never
mentions the frame.

The obstruction is uniformity. A lifting rule earns the name ``rule'' by being
indifferent to which proposition it lifts. That indifference is exactly what
prevents it from delivering a particular new fact --- if it delivers $\varphi$
it must, under substitution, deliver every substitution instance of $\varphi$,
and those are jointly unsatisfiable whenever $\varphi$ is not already forced.

The regress is harmless, then, for a reason more specific than consistency. It
is harmless because it is inert, and it is inert because it is uniform.

\subsection{Where transcendence can still live}

The theorems are not a proof that context logic is useless; they delimit where
its content can come from. Three doors remain open, corresponding to the three
corners.

\emph{Through (U).} Give up uniqueness. Multiple extensions are not a defect if
one is prepared to reason about all of them, or to impose priorities
extra-logically. This is the standard nonmonotonic bargain, and the price is
that ``what follows'' becomes route-dependent: the same fact may arrive
guaranteed by one path and contingent by another.

\emph{Through (P).} Give up uniformity, deliberately and visibly. Domain axioms
that name particular propositions or particular contexts are perfectly
legitimate; they are simply assertions, and should be labelled as such rather
than presented as logical machinery. Much of what has been called a lifting
axiom in the literature is of this kind.

\emph{Through the base.} Enrich $T$ so that it is not symmetric. Real
applications do this constantly, and Theorem \ref{thm:trichotomy} then does not
apply. What our results show is that the novelty in such a system comes from the
asymmetry of $T$, not from the tower.

\subsection{Hypotheses a referee will press}

\begin{remark}[Context-freeness of $T$]\label{rem:contextfree}
If $T \not\subseteq L$ the substitution argument of Theorem \ref{thm:mono}
breaks at the step where $s\psi_i$ is shown tautologous, and the theorem is
false: take $T = \{\ist_{c_1}(\theta)\}$ with an exit rule. We regard this as
the correct scope rather than a limitation --- the question is what the tower
contributes, and a base theory already speaking of contexts has prejudged it.
\end{remark}

\begin{remark}[Symmetry of $T$]
Symmetry is the one hypothesis of Theorem \ref{thm:trichotomy} that is not
forced by the subject matter, and it is where further work should go. It is used
only in Step 2, and only to move an arbitrary model of $T$ to a fixed one. Any
condition guaranteeing that $\Cn(T)$ is the largest $\Aut(L)_T$-invariant
consistent deductively closed set would serve. We do not know the optimal
hypothesis.
\end{remark}

\begin{remark}[Propositional setting]
We work propositionally. The substitution argument should survive the passage to
first-order, with substitutions of formulas for predicate letters in place of
substitutions for atoms, but the automorphism group in
Theorem \ref{thm:trichotomy} becomes more delicate and we have not checked it.
\end{remark}

\subsection{Relation to the machine-checked development}\label{sec:formal}

All three theorems above are now formalized in Lean 4 against Mathlib, in
\texttt{Development/\allowbreak Uniformity.lean},
\texttt{Development/\allowbreak Rigidity.lean} and
\texttt{Development/\allowbreak Trichotomy.lean}. The formalization is semantic throughout:
$\Cn$ is truth in every valuation rather than derivability. Four things it
established are worth recording here, because three of them revise claims made
above.

\emph{Theorem \ref{thm:mono} needs neither a proof calculus nor compactness.}
The proof in Section 3 extracts a finite derivation and substitutes into it. The
machine-checked proof instead composes valuations: substitution and evaluation
commute, in the sense that $v \models s\varphi$ exactly when $s^{*}v \models
\varphi$ for the pulled-back valuation $s^{*}v$, and the theorem is eight lines
from there. The frame is absent from the statement and the proof alike --- no
$\out$ map occurs --- which is stronger than the remark following the theorem.

\emph{Theorem \ref{thm:rigidity} is proved for frames whose map is a
bijection.} Surjectivity is used twice, and neither use is removable from the
present argument. It is needed for $\Cn$ to commute with $\sigma^{-1}$, since
every model of $\sigma^{-1}(X)$ must be the image of a model of $X$; and it is
needed in the stage construction, where a rule of $S_F$ whose consequent lies
outside the image of $\sigma$ has no preimage rule to fire. The second is
exactly what the clause ``groundedness is inherited stagewise by the same
argument'' passes over. We do not know whether the theorem holds for a merely
injective frame map. By Remark \ref{rem:definable} the bi-infinite frame, where
the shift is a bijection, is the case carrying the content.

\emph{The compactness step in Theorem \ref{thm:trichotomy} does not exist.}
Step 2 above collects the orbit $\{g\varphi\}$, observes that
$T \cup \{g\varphi\}_g$ is unsatisfiable, and appeals to compactness to convert
that into inconsistency. Semantically there is nothing to convert: a theory with
no model is inconsistent by definition. The formalized argument does not even
form the orbit. It takes the model $v$ of the extension --- which is a model of
$T$, since the extension contains $\Cn(T)$ --- and the single $g$ carrying the
counter-model $M$ to $v$, and derives that $g\varphi$ is both in the extension
and false in $v$. The claim that this is the only place where the infinite index
set does any work should be withdrawn: it does none.

\emph{Step 1 of Theorem \ref{thm:trichotomy} and Theorem \ref{thm:rigidity} are
one lemma.} Both assert that an automorphism preserving $\Cn(T)$ and preserving
the schema carries extensions to extensions, so that uniqueness forces
invariance. The development states this once, for an arbitrary automorphism of
the context language realised on valuations, and obtains rigidity by
instantiating it at the shift and Step 1 by instantiating it at an atom
automorphism.

Two further points. Formalizing Step 1 is what forced the signed-atom syntax
mentioned after Definition \ref{def:aut}: the preimage of the schema under an
element of $\Aut(L)$ must be the schema \emph{as a set of rules}, and that fails
if the group acts only up to equivalence. And the assertion in Definition
\ref{def:extension} that a monotone schema has a unique fixed point is proved
rather than assumed, as is the consistency of the copy schema, so that none of
the three theorems is about an empty class of schemas.

The corners are occupied. The copy schema is shown to lie in (C) by the
trichotomy itself rather than by direct computation; the earlier development
supplies (U) and (P) through explicit multi-extension and route-sensitivity
constructions.

\end{document}
